% !TeX spellcheck = it_IT

\documentclass[]{beamer}

\usepackage[utf8]{inputenc}
\usepackage[italian]{babel}
\usepackage{multicol}
\usetheme{metropolis}

\usepackage{mathtools}
\DeclarePairedDelimiter{\ceil}{\lceil}{\rceil}


\title{Informatica B - Esercitazione 0}
\subtitle{Codifica binaria - Numeri naturali}

\author{Stefano Cereda\\
		stefano1.cereda@mail.polimi.it
	}
\date{10/10/2017}
\institute[PoliMi]{\vspace{0.5cm}\centering Politecnico di Milano \\ \vspace{0.2cm}
	\includegraphics[width=0.2\linewidth]{./logopolimi}}
%\logo{\includegraphics[width=15mm]{./logopolimi}}

%\setbeamercovered{invisible}

\makeindex

\begin{document}
\begin{frame}
	\maketitle
\end{frame}

\section{Recap e Introduzione}
\begin{frame}{Recap}
\begin{itemize}[<+->]
	\item Combinazioni possibili e numero di bit richiesti: \\$2^N$ (da 0 a $2^N-1$) e $\ceil{log_2M}$
	\item Codifica binaria pura: \\$1011_2 = 1+2+8=11_{10}$ e divisione ripetuta
	\item Somme in binario puro: \\$0100+111 = 1011$
	\item ASCII: \\ `A'+1 = `B'
\end{itemize}
\end{frame}

\begin{frame}{Introduzione}
\begin{itemize}[<+->]
	\item Numeri interi: codifica modulo e segno e codifica in complemento alla base
	\item Non faremo i numeri reali	
\end{itemize}
\end{frame}

\section{Numeri interi}
\begin{frame}{Modulo e segno}
Nella codifica in modulo e segno si utilizza il bit più significativo (\emph{msb}) per indicare il segno dei numeri:

$0101_2 = +101_2 = +(1+4)_{10} = +5_{10}$\\
$1101_2 = -101_2 = -(1+4)_{10} = -5_{10}$

\pause

Con N bit useremo un bit per il segno, i rimanenti per il modulo. Possiamo quindi rappresentare numeri nell'intervallo $-2^{N-1}+1 \leq x \leq +2^{N-1}-1$

\pause

Abbiamo infatti due possibili codifiche per il numero zero:
$+0 = +0000 = 00000$\\
$-0 = -0000 = 10000$\\
($N=5$)
\end{frame}

\begin{frame}{Complemento alla base}
La precedente codifica ``spreca'' un valore, con N bit vorremo infatti poter rappresentare $2^N$ valori, non $2^N-1$.

La rappresentazione in complemento alla base rappresenta, con N bit, i valori da $-2^{N-1}$ a $2^{N-1}-1$

\pause

La rappresentazione in complemento a 2, su N bit, di un numero $x_{10}$ è definita come:
\begin{itemize}
	\item $x_2$ se $x \geq 0$ (msb = 0)
	\item $(2^N-|x|)_2$ se $x < 0$ (msb = 1, dimostrare)
\end{itemize}

\pause

Abbiamo quindi una sola rappresentazione per il numero 0.

Attenzione: il $msb$ indica il segno solo se siamo in base 2, non è un bit di segno in senso stretto.
\end{frame}

\begin{frame}{Esempio cpl2 su 3 bit}
$N=3 \rightarrow -2^2 \leq x \leq 2^2-1 \rightarrow -4 \leq x \leq 3 \ \ \ M=2^3=8$
\pause

\centering
\begin{tabular}{c|c|l}
	10	&	cpl2	&	da 10 a cpl2 \\
	\hline
	+3&	011	&	1+2\\
	+2& 010	&	2\\
	+1& 001	&	1\\
	0 & 000	&	0\\
	-1& 111	&	$2^3-|-1| = 8-1 = 7_{10} = 111_2$\\
	-2& 110	&	$8-2 = 6_{10} = 110_2$\\
	-3& 101	&	$8-3 = 5$\\
	-4& 100	&	$8-4 = 4$\\
\end{tabular}

Notare che msb \alert{non è} il segno: cambiandolo non si ottiene il numero opposto.
\end{frame}

\begin{frame}{Proprietà cpl2}
Rappresentiamo $-x$ con $N=4$ bit:\\ $-x \rightarrow (2^N-|-x|)_2 = 10000 -x = 1111-x+1$

\pause
$1111-x$ possiamo ottenerlo semplicemente invertendo tutti i bit di $x$:
\begin{tabular}{c|c}
	$1111$ & - \\
	$0101$ & = \\
	\hline
	$1010$& \\
\end{tabular}

\pause
Adottando la codifica in complemento a 2 il calcolatore necessita solo di un circuito di addizione (più quello banale di complementazione) per effettuare sia somme che sottrazioni!
\end{frame}

\begin{frame}{Metodi di conversione da base 10 a cpl2}
Per convertire un numero negativo da base 10 a cpl2 abbiamo tre metodi:
\begin{enumerate}[<+->]
	\item Utilizzare la definizione: $(2^N-|x|)_2$:\\
	$N=3; -2 \rightarrow (2^N-2)_2 = (6)_2 \rightarrow 110$
	\item Convertire l'opposto in base 2, invertire i bit e sommare 1:\\
	$N=3; -2 \rightarrow +2_{10} = 010_2 \rightarrow 101 \rightarrow 110$
	\item Convertire l'opposto in base due, ricopiarlo da $lsb$ verso $msb$ fino al primo 1 (compreso), copiare i restanti bit complementati:\\
	$N=4; -6 \rightarrow +6_{10} = 0110_2 \rightarrow 1010$
\end{enumerate}

\pause
Controlliamo l'ultimo risultato utilizzando la definizione di cpl2: $1010_2 = 2^N - |x| = 2^N+x \rightarrow x=1010_2-2^N = 10-16 = -6$
\end{frame}

\begin{frame}{Operazioni algebriche in cpl2}
Dati due numeri in base 10 da sommare algebricamente, dobbiamo innanzitutto controllare il numero di bit necessari:
\begin{itemize}
	\item Se N è assegnato, dobbiamo verificare che sia sufficiente. ($-2^{N-1} \leq x \leq 2^{N-1}-1$)
	\item Altrimenti dobbiamo calcolare il valore minimo capace di rappresentare \alert{entrambi} i valori.
\end{itemize}
\end{frame}

\begin{frame}{Somma senza riporto ne overflow}
$(-7) + (+2) = -5$ con $N=4$

\pause
$-2^{N-1} = -8 \wedge 2^{N-1}-1 = +7 \rightarrow$ i bit sono sufficienti

\pause
$-7 \rightarrow 7_{10} = 0111_2 \rightarrow 1001_{cpl2}$\\
$+2 \rightarrow 2_{10} = 0010_2 \rightarrow 0010_{cpl2}$

\pause
\begin{tabular}{c|c}
	$1001$ & + \\
	$0010$ & = \\
	\hline
	$1011$& \\
\end{tabular}

\pause
$1011 < 0 \rightarrow 0101 = 5 \rightarrow -5$
\end{frame}

\begin{frame}{Somma con riporto ma senza overflow}
$(+7) + (-2) = +5$ con $N=4$

$-2^{N-1} = -8 \wedge 2^{N-1}-1 = +7 \rightarrow$ i bit sono sufficienti

$+7 \rightarrow 7_{10} = 0111_2 \rightarrow 0111_{cpl2}$\\
$-2 \rightarrow 2_{10} = 0010_2 \rightarrow 1110_{cpl2}$

\begin{tabular}{c|c}
	$0111$ & + \\
	$1110$ & = \\
	\hline
	\hskip-0.15cm\alert{1}$0101$& \\
\end{tabular}

Ignoriamo il riporto:\\
$0101 > 0 \rightarrow 0101 = 5 \rightarrow +5$
Il risultato è corretto
\end{frame}

\begin{frame}{Somma senza riporto ma con overflow}
$(+7) + (+2) = +9$ con $N=4$

$-2^{N-1} = -8 \wedge 2^{N-1}-1 = +7 \rightarrow$ i bit sono sufficienti per gli operandi, ma avremo una situazione di overflow

$+7 \rightarrow 7_{10} = 0111_2 \rightarrow 0111_{cpl2}$\\
$+2 \rightarrow 2_{10} = 0010_2 \rightarrow 0010_{cpl2}$

\begin{tabular}{c|c}
	$0111$ & + \\
	$0010$ & = \\
	\hline
	$1001$& \\
\end{tabular}

$1001 < 0 \rightarrow 0111 = 7 \rightarrow -7$

Non abbiamo riporto, cosa ci indica la presenza di overflow?
\end{frame}

\begin{frame}{Somma con riporto e overflow}
$(-7) + (-2) = -9$ con $N=4$

$-2^{N-1} = -8 \wedge 2^{N-1}-1 = +7 \rightarrow$ i bit sono sufficienti per gli operandi, ma avremo una situazione di overflow

$-7 \rightarrow 7_{10} = 0111_2 \rightarrow 1001_{cpl2}$\\
$-2 \rightarrow 2_{10} = 0010_2 \rightarrow 1110_{cpl2}$

\begin{tabular}{c|c}
	$1001$ & + \\
	$1110$ & = \\
	\hline
	\hskip-0.15cm\alert{1}$0111$& \\
\end{tabular}

Ignoriamo il riporto:
$0111 > 0 \rightarrow 0111 = 7 \rightarrow +7$

Operandi concordi ma risultato discorde indicano che siamo in una condizione di overflow.

(Operandi discordi non daranno mai overflow)
\end{frame}
\end{document}
